\section{Conclusion}
\label{sec:conclusion}

We presented \textsc{Intro-Specter}, a three-layer self-correction
framework that formalises agent self-correction as abductive
Bayesian inference over an explicit assumption graph. Where
existing methods either critique flat trajectories
(Self-Refine, Reflexion) or search forward over reasoning
branches (Tree of Thoughts), \textsc{Intro-Specter} jointly
(i) represents profile-grounded assumptions as nodes with
provenance, (ii) infers a posterior over candidate faulty
assumptions after a violation is detected, and (iii) selectively
repairs only the downstream subgraph while preserving the valid
prefix.

Across \textbf{eight reconstruction benchmarks plus four
real-data fidelity anchors, eight LLMs spanning four training
families and 7B--671B parameters, and seven self-correction
baselines compared head-to-head}, \textsc{Intro-Specter} delivers
\textbf{18 Holm-Bonferroni-significant wins against Direct}
spanning all seven major reconstruction benchmarks plus the
\textsc{Real-HotpotQA} real-data anchor, peaks at \textbf{$+33.3$
percentage points on \textsc{Profile-MuSiQue} $\times$ Gemini 2.5
Flash} ($p = 10^{-5}$), achieves \textbf{100\% top-1 attribution
accuracy and 100\% repair success at zero token overhead} in the
controlled synthetic setting (matching both Oracle-Detector and
Oracle-Repair upper bounds), and \textbf{recovers about half the
cross-family model-quality gap} on profile-conditioned QA.

The detection-only baseline sits \emph{exactly} at Direct on
every \textsc{Real-HotpotQA} cell, providing the cleanest
empirical evidence in this work that detection alone is
insufficient: the joint engineering of detection, attribution,
and selective repair is what unlocks the gains.

\textbf{When the method does not win.} Reflexion remains a strong
co-baseline. It strictly beats \textsc{Intro-Specter} at
$p < 0.05$ on 2 of 28 \textsc{Real-*} cells --- specifically
\textsc{Real-TruthfulQA} $\times$ DeepSeek V3 (broken-prefix
regime: nothing salvageable for selective repair) and
\textsc{Real-TravelPlanner} $\times$ DeepSeek V3 (long-horizon
diffuse-fault regime). On ceiling-effect models (Llama 3.3 70B
near-perfect on factual QA) no method moves the needle.
\textsc{Intro-Specter} is the right tool for the deployment
regime that the paper is theoretically about: weaker open-weight
models on profile-grounded reasoning tasks with deep but
recoverable trajectories. In other regimes the two methods are
complementary, and a deployment system would route between them
based on the trajectory's recoverability signal.

\textbf{Future work.} Three directions stand out. First, the
val-split-tuned $\tau$-abstention calibration sweep at the time
of writing covers four \textsc{Profile-MuSiQue} cells; extending
to all 18 IS-winning cells would give a fully held-out
calibrated-abstention claim. Second, ToT's competitive
performance on \textsc{Profile-MuSiQue} $\times$ Mistral Nemo
(87.5\% vs IS's 88.3\%) suggests an interesting hybrid:
running ToT \emph{forward} from the IS-selected repair node could
combine the strengths of both methods. Third, the
profile-update mechanism --- handling users whose preferences
drift over time --- is a natural extension we have not
addressed.

\textbf{Reproducibility.} Code, configs, aggregated CSVs, and
the auto-generated booktabs LaTeX are released at
\url{https://github.com/Jihan-5/Neurlips-Intro-Specter}. All
LLM calls are cached deterministically; the entire matrix can
be re-aggregated from JSONLs in $<5$ minutes. We hope the
reconstruction methodology, the unified profile + fault
injection harness, and the eight-baseline comparison pipeline
are useful artifacts for future work on agent self-correction.
